\documentclass[12pt,a4paper]{article}

% Packages
\usepackage[a4paper, margin=2.5cm]{geometry}
\usepackage{setspace}
\usepackage{titlesec}
\usepackage{parskip}
\usepackage{hyperref}
\usepackage{fontenc}
\usepackage{inputenc}
\usepackage{lmodern}
\usepackage{microtype}
\usepackage{booktabs}
\usepackage{enumitem}
\usepackage{fancyhdr}
\usepackage{graphicx}
\usepackage{xcolor}

% Page style
\pagestyle{fancy}
\fancyhf{}
\rhead{\textit{Research Methodology}}
\lhead{\textit{IDP -- Soil Property Analysis}}
\cfoot{\thepage}
\renewcommand{\headrulewidth}{0.4pt}

% Section formatting
\titleformat{\section}{\large\bfseries}{}{0em}{}[\titlerule]
\titleformat{\subsection}{\normalsize\bfseries}{}{0em}{}

% Spacing
\onehalfspacing

% Hyperlink colours
\hypersetup{
    colorlinks=true,
    linkcolor=black,
    urlcolor=blue,
    citecolor=black
}

% -------------------------------------------------------
\begin{document}

% Title Page
\begin{titlepage}
    \centering
    \vspace*{2cm}

    {\LARGE\bfseries Web-Based System for Soil Property Analysis\\[0.4em]
    Using Machine Learning Techniques\par}

    \vspace{1.5cm}
    {\large\itshape Research Methodology\par}

    \vspace{0.5cm}
    {\normalsize Preliminary Methodology Draft for IDP Mentor Review\par}

    \vspace{2cm}
    \rule{0.6\textwidth}{0.4pt}

    \vspace{1cm}
    {\normalsize
    \textbf{Author:} Sampurna Chandra\\[0.3em]
    \textbf{GitHub:} \texttt{sampurnachandra97-max}\\[0.3em]
    \textbf{Repository:} \href{https://github.com/sampurnachandra97-max/idp-2026}%
        {\texttt{sampurnachandra97-max/idp-2026}}
    }

    \vspace{1.5cm}
    {\small\itshape Interdisciplinary Project (IDP) -- 2026\par}

    \vfill
    {\small \today}
\end{titlepage}

% -------------------------------------------------------
\tableofcontents
\newpage

% -------------------------------------------------------
\section{Research Approach}

The proposed project is inspired by the research paper \textit{``Estimation of
Soil Properties Using Machine Learning Techniques to Improve Hydrological
Modelling in a Semiarid Environment: Campo de Cartagena (Spain).''} The main
purpose of this project is to explore how machine learning techniques can be
used to study soil properties and how this information can be presented through
a web-based system.

The original research focuses on predicting soil organic carbon, clay, silt, and
sand using machine learning algorithms.  Our project takes this research as a
reference for developing a user-friendly website that can support the
understanding and analysis of soil-related information.

% -------------------------------------------------------
\section{Study of Existing Research}

The initial stage involves studying the existing research paper to understand
the problem, the datasets used, the machine learning algorithms applied, and the
results obtained. The paper compares four algorithms for soil-property
estimation:

\begin{itemize}[noitemsep]
    \item \textbf{Random Forest (RF)}
    \item \textbf{Support Vector Machines (SVM)}
    \item \textbf{XGBoost}
    \item \textbf{Multilayer Perceptron (MLP)}
\end{itemize}

The study also explains how predicted soil properties can be used to improve
hydrological modelling.  This literature study provides the theoretical
foundation for understanding the problem and planning the proposed system.

% -------------------------------------------------------
\section{Proposed System Design}

The project involves developing a web-based interface through which soil-related
information can be presented in an organised and accessible manner.

The system is designed with a focus on:
\begin{itemize}[noitemsep]
    \item Simple navigation and clear presentation of information.
    \item An interface that users can understand without requiring extensive
          technical knowledge.
    \item Responsive layout suitable for desktop and mobile screens.
\end{itemize}

The website forms the main platform for the project.  Its features and
functionalities are documented according to the final implementation.

% -------------------------------------------------------
\section{Website Development and Technologies Used}

The website has been developed using \textbf{Next.js}, a React-based framework
for building web applications.  \textbf{Google Antigravity}, an AI-assisted
development environment, was used during the development process to support the
creation and implementation of the website.

The development process involved:
\begin{enumerate}[noitemsep]
    \item Building the website interface and organising its components.
    \item Implementing the required pages and functionalities.
    \item Ensuring component-based structure for easier maintenance.
\end{enumerate}

The use of Next.js supports a component-based architecture, making the
application easier to organise and maintain.

% -------------------------------------------------------
\section{Machine Learning Approach}

The reference paper provides a foundation for understanding how machine learning
can be applied to soil-property prediction.  The algorithms considered in the
literature study include Random Forest, Support Vector Machines, XGBoost, and
Multilayer Perceptron.

These algorithms are studied to understand their suitability for predicting soil
properties from environmental data.  The implementation of a machine learning
prediction module, if included in the final project, will require:

\begin{itemize}[noitemsep]
    \item A suitable and well-structured dataset.
    \item Appropriate data preprocessing methods.
    \item Model training and hyperparameter tuning.
    \item Independent performance evaluation using standard metrics.
\end{itemize}

% -------------------------------------------------------
\section{System Testing and Evaluation}

The website will be evaluated based on its functionality, usability, and ability
to present information clearly.  Testing will involve:

\begin{itemize}[noitemsep]
    \item Verifying that all implemented features work correctly.
    \item Checking that navigation is smooth and intuitive.
    \item Ensuring the interface behaves properly on different screen sizes.
\end{itemize}

If a machine learning prediction model is integrated into the system, its
performance will be evaluated separately using appropriate metrics such as
\textbf{RMSE} (Root Mean Square Error) and \textbf{R\textsuperscript{2}}
(Coefficient of Determination).

% -------------------------------------------------------
\section{Expected Outcome}

The expected outcome of the project is a functional web-based system that
provides an accessible platform for working with soil-related information.

The project aims to connect the theoretical concepts discussed in the research
paper with a practical software application.  Future development may include:

\begin{itemize}[noitemsep]
    \item Integrating trained machine learning models for real-time
          soil-property estimation.
    \item Incorporating environmental datasets to support analysis.
    \item Extending the platform with data visualisation dashboards.
\end{itemize}

\vspace{1.5cm}
\hrule
\vspace{0.5cm}
{\small\itshape This document is a preliminary draft prepared for IDP mentor
review and is subject to revision.}

\end{document}
